% Part: lambda-calculus
% Chapter: lambda-definability
% Section: lists

\documentclass[../../../include/open-logic-section]{subfiles}

\begin{document}

\olfileid[tr]{lam}{ldf}{lst}
\olsection{Listeler}

$[M_1, M_2, \ldots, M_n]$ listesi şöyle tanımlanabilir:
\[
  [M_1, M_2, \ldots, M_n] \ident \lambd[sz][s M_1 (s M_2 (\ldots (s M_n z)))]
\]
Gayriresmî olarak, bir listeyi temsil eden terim listenin ``biriktirici''
fonksiyonudur: Bu fonksiyon iki terim kabul eder. Birincisi, biriktirilmiş
değerle yeni bir elemanı alıp yeni biriktirilmiş değeri döndüren fonksiyondur;
ikincisi başlangıçtaki biriktirilmiş değerdir. Elemanları birer birer
biriktirir ve sonunda sonucu döndürür.

\begin{digress}
  Fonksiyonel programlamaya aşinaysanız bu tanımın \emph{sağ katlama}ya
  benzediğini göreceksiniz: Bir liste, içerdiği elemanların sağ katlama
  fonksiyonu olarak tanımlanır.
\end{digress}

Bu kodlamaya dayanarak bazı yararlı fonksiyonları kolayca tanımlayabiliriz:
\begin{align*}
  \fn{Sum} & \ident
  \lambd[l][l \, (\lambd[xa][\fn{Add} \, x \, a])\,  \num{0}]\\
  \fn{Len} & \ident
  \lambd[l][l \, (\lambd[xa][\fn{Add} \, a \, \num{1}]) \num{0}]
\end{align*}

$\fn{Sum}$, Church sayılarından oluşan bir listenin toplamını hesaplar.
Bunu, başlangıç değeri $\num{0}$ olan bir liste biriktirmesi yaparak ve her
eleman için yeni değer olarak $\fn{Add} \, x\, a$ değerini döndürerek yapar
(burada $x$ geçerli eleman, $a$ ise biriktirilmiş değerdir). Sonuç bütün
elemanların toplamıdır.

\begin{prob}
  $\fn{Len}$ ne yapar? Açıklayın.
\end{prob}

\begin{prob}
  Bir listeyi tersine çeviren bir fonksiyon tanımlayın.
\end{prob}
\end{document}

